\documentclass[11pt]{article}

% --- Margins ---
\usepackage[margin=1in]{geometry}

% --- Math ---
\usepackage{amsmath, amssymb, amsthm, mathtools}
\usepackage{bm}            % bold math symbols (\bm)

% --- Quality of life ---
\usepackage{enumitem}      % list customization
\usepackage{microtype}     % better justification
\usepackage[colorlinks=true, allcolors=blue]{hyperref}
\usepackage{cleveref}      % \cref / \Cref smart references

% --- Theorem environments ---
\theoremstyle{plain}
\newtheorem{theorem}{Theorem}
\newtheorem{lemma}{Lemma}
\newtheorem{proposition}{Proposition}
\newtheorem{corollary}{Corollary}
\theoremstyle{definition}
\newtheorem{definition}{Definition}
\newtheorem{assumption}{Assumption}
\theoremstyle{remark}
\newtheorem{remark}{Remark}

% --- Macros ---
\newcommand{\R}{\mathbb{R}}
\newcommand{\T}{^{\mathsf{T}}}                       % transpose
\newcommand{\inv}{^{-1}}
\newcommand{\grad}{\nabla}
\newcommand{\half}{\tfrac{1}{2}}
\newcommand{\diag}{\operatorname{diag}}
\newcommand{\tr}{\operatorname{tr}}
\newcommand{\rank}{\operatorname{rank}}
\newcommand{\st}{\;\text{s.t.}\;}
\DeclareMathOperator*{\argmin}{arg\,min}
\DeclareMathOperator*{\argmax}{arg\,max}
\DeclareMathOperator*{\minimize}{minimize}
\DeclareMathOperator{\inertia}{In}
\newcommand{\norm}[1]{\left\lVert #1 \right\rVert}
\newcommand{\abs}[1]{\left\lvert #1 \right\rvert}
\newcommand{\inner}[2]{\left\langle #1, #2 \right\rangle}
\newcommand{\set}[1]{\left\{ #1 \right\}}
\newcommand{\pd}[2]{\frac{\partial #1}{\partial #2}}  % partial derivative
% Complementarity shorthand: a \perp b
\newcommand{\perpcomp}{\;\perp\;}

\title{Marble Solver Derivation}
\author{Robotic Exploration Lab}
\date{\today}

\begin{document}
\maketitle

\section{Problem Setup}

We seek to solve the following Quadratic Program with Complementarity Constraints (QPCC):

\begin{equation}\label{qpcc}
\begin{aligned}
\minimize_{x \in \R^n} \quad & \frac12 x^\top Q x + g^\top x + c \\
\st \quad & Ax + b = 0 \\
          & Gx + h \geq 0 \\
          & 0 \leq Lx + l \perp Rx + r \geq 0
\end{aligned}
\end{equation}
for $Q \succeq 0$, $A \in \R^{m \times n}$, $G \in \R^{p \times n}$, $L, R \in \R^{c \times n}$, $b \in \R^m$, $h \in \R^p$, and $l, r \in \R^c$.

\section{Derivation}

We introduce the $\textit{retraction map}$ $b_\kappa(x)$ which satisfies:
\begin{equation}\label{retraction_map}
    b_\kappa(x) \cdot b_\kappa(-x) = \kappa
\end{equation}
Using \eqref{retraction_map}, we can relax the complementarity constraints in \eqref{qpcc} as:
\begin{equation}\label{qpcc_relax}
\begin{aligned}
\minimize_{x \in \R^n, \sigma\in \R^c, d \in \R^p} \quad & \frac12 x^\top Q x + g^\top x + c - \kappa\bm{1}^T \log(d) \\
\st \quad & Ax + b = 0 \\
          & Gx + h = d \\
          & Lx + l = b_\kappa(\sigma) \\
          & Rx + r = b_\kappa(-\sigma)
\end{aligned}
\end{equation}
We then introduce the Augmented Lagrangian for \eqref{qpcc_relax} as:
\begin{equation}\label{aug_lag}
\begin{aligned}
\mathcal{L}_\rho(x, \sigma, d; \alpha, \mu, \gamma_L, \gamma_R) = & \frac12 x^\top Q x + g^\top x + c \\
& + \mu^\top (Gx + h - d) - \kappa\bm{1}^T \log(d)\\
& + \alpha^\top (Ax + b) + \frac\rho2 \norm{Ax + b}^2 \\
& + \gamma_L^\top (Lx + l - b_\kappa(\sigma)) + \frac\rho2 \norm{Lx + l - b_\kappa(\sigma)}^2 \\
& + \gamma_R^\top (Rx + r - b_\kappa(-\sigma)) + \frac\rho2 \norm{Rx + r - b_\kappa(-\sigma)}^2
\end{aligned}
\end{equation}
Using an Augmented-Lagrangian approach, we can solve \eqref{qpcc_relax} by iteratively minimizing \eqref{aug_lag} 
over primal variables $x, \sigma, d$.
We update AL multipliers $\alpha, \gamma_L, \gamma_R$ in the outer loop, while $\mu$ is solved as an
IPM dual variable in the inner loop.

We first derive stationarity with respect to the original primal variable $x$. Define the following 
shorthands:
\begin{align}
    r_A &:= Ax + b \\
    r_G &:= Gx + h - d \\
    r_L &:= Lx + l - b_\kappa(\sigma) \\
    r_R &:= Rx + r - b_\kappa(-\sigma)
\end{align}
We therefore have the following stationarity condition:
\begin{align}
\pd{\mathcal{L}_\rho}{x} &= Qx + g + A^\top \alpha + \rho A^\top r_A + G^\top \mu + L^\top \gamma_L + \rho L^\top r_L + R^\top \gamma_R + \rho R^\top r_R = 0 \\
&= Qx + g + A^\top \lambda + G^\top \mu + L^\top \tau_L + R^\top \tau_R = 0
\end{align}
We introduce the primal-dual shifted multiplier variables:
\begin{align}
    \lambda &= \alpha + \rho r_A = \alpha + \rho (Ax + b) \\
    \tau_L &= \gamma_L + \rho r_L = \gamma_L + \rho (Lx + l - b_\kappa(\sigma)) \\
    \tau_R &= \gamma_R + \rho r_R = \gamma_R + \rho (Rx + r - b_\kappa(-\sigma))
\end{align}
For variable $d$, we have the following stationarity condition:
\begin{align}
\pd{\mathcal{L}_\rho}{d} &= -\mu - \kappa d^{-1} = 0 \\
&\Longleftrightarrow \mu \odot d = -\kappa \bm{1}
\end{align}

We can satisfy this relaxed complementarity condition by again using \eqref{retraction_map} as:
\begin{equation}
    d = b_\kappa(v), \quad \mu = -b_\kappa(-v), \quad v \in \R^p
\end{equation}

And finally for variable $\sigma$, we have the following stationarity condition:
\begin{align}
\pd{\mathcal{L}_\rho}{\sigma} &= -b'_\kappa(\sigma) (\gamma_L + \rho r_L) + b'_\kappa(-\sigma) (\gamma_R + \rho r_R) = 0 \\
&= -b'_\kappa(\sigma) \tau_L + b'_\kappa(-\sigma) \tau_R = 0
\end{align}

In all, we have the following KKT conditions for \eqref{qpcc_relax}:
\begin{align}
    r_x &:= Qx + g + A^\top \lambda + G^\top \mu + L^\top \tau_L + R^\top \tau_R &&= 0 \\
    r_\mu &:= \mu + b_\kappa(-v) &&= 0 \\
    r_\sigma &:= -b'_\kappa(\sigma) \tau_L + b'_\kappa(-\sigma) \tau_R &&= 0 \\
    r_\lambda &:= Ax + b - \rho^{-1} (\lambda - \alpha) &&= 0 \\
    r_i &:= Gx + h - b_\kappa(v) &&= 0 \\
    r_{\tau_L} &:= Lx + l - b_\kappa(\sigma) - \rho^{-1} (\tau_L - \gamma_L) &&= 0 \\
    r_{\tau_R} &:= Rx + r - b_\kappa(-\sigma) - \rho^{-1} (\tau_R - \gamma_R) &&= 0
\end{align}
We now write the Newton system for the KKT conditions described above.
We note that
$b''_\kappa(\sigma) = b''_\kappa(-\sigma)$ since $b''_\kappa(x)$ is an even function.
\begin{equation}
\begin{bmatrix}
    Q & 0 & 0 & A^\top & G^\top & L^\top & R^\top \\
    0 & -b'_\kappa(-v) & 0 & 0 & I & 0 & 0 \\
    0 & 0 & -b''_\kappa(\sigma) \odot (\tau_L + \tau_R) & 0 & 0 & -b'_\kappa(\sigma) & b'_\kappa(-\sigma) \\
    A & 0 & 0 & -\rho^{-1} I & 0 & 0 & 0 \\
    G & -b'_\kappa(v) & 0 & 0 & 0 & 0 & 0 \\
    L & 0 & -b'_\kappa(\sigma) & 0 & 0 & -\rho^{-1} I & 0 \\
    R & 0 & b'_\kappa(-\sigma) & 0 & 0 & 0 & -\rho^{-1} I
\end{bmatrix}
\begin{bmatrix}
    \Delta x \\ \Delta v \\ \Delta \sigma \\ \Delta \lambda \\ \Delta \mu \\ \Delta \tau_L \\ \Delta \tau_R
\end{bmatrix}
=
-\begin{bmatrix}
    r_x \\ r_\mu \\ r_\sigma \\ r_\lambda \\ r_i \\ r_{\tau_L} \\ r_{\tau_R}
\end{bmatrix}
\end{equation}

Using the second equation $-b'_\kappa(-v) \Delta v + \Delta \mu = -r_\mu$, we solve for $\Delta \mu$ and get:
\[
  Q \Delta x + A^\top \Delta \lambda + G^\top b'_\kappa(-v) \Delta v + L^\top \Delta \tau_L + R^\top \Delta \tau_R = -r_x + G^\top r_\mu
\]
From this, we subtract $G^\top$ times the fifth equation $G \Delta x - b'_\kappa(v) \Delta v = -r_i$ to get:
\[
    (Q - G^\top G) \Delta x + A^\top \Delta \lambda + G^\top \Delta v+ L^\top \Delta \tau_L + R^\top \Delta \tau_R = -r_x + G^\top (r_\mu + r_i)
\]
where we've used the identity $b'_\kappa(-v) + b'_\kappa(v) = I$. The system is now symmetric, and written as:
\begin{equation}
\begin{bmatrix}
    Q - G^\top G & G^\top & 0 & A^\top & L^\top & R^\top \\
    G & -b'_\kappa(v) & 0 & 0 & 0 & 0 \\
    0 & 0 & -b''_\kappa(\sigma) \odot (\tau_L + \tau_R) & 0 & -b'_\kappa(\sigma) & b'_\kappa(-\sigma) \\
    A & 0 & 0 & -\rho^{-1} I & 0 & 0 \\
    L & 0 & -b'_\kappa(\sigma) & 0 & -\rho^{-1} I & 0 \\
    R & 0 & b'_\kappa(-\sigma) & 0 & 0 & -\rho^{-1} I
\end{bmatrix}
\begin{bmatrix}
    \Delta x \\ \Delta v \\ \Delta \sigma \\ \Delta \lambda \\ \Delta \tau_L \\ \Delta \tau_R
\end{bmatrix}
=
-\begin{bmatrix}
    r_x - G^\top (r_\mu + r_i) \\ r_\sigma \\ r_\lambda \\ r_i \\ r_{\tau_L} \\ r_{\tau_R}
\end{bmatrix}
\end{equation}








\newpage
\section{Derivation}

We introduce the $\textit{retraction map}$ $b_\kappa(x)$ which satisfies:
\begin{equation}\label{retraction_map}
    b_\kappa(x) \cdot b_\kappa(-x) = \kappa
\end{equation}
Using \eqref{retraction_map}, we can relax the complementarity constraints in \eqref{qpcc} as:
\begin{equation}\label{qpcc_relax}
\begin{aligned}
\minimize_{x \in \R^n, \sigma\in \R^c, d \in \R^p} \quad & \frac12 x^\top Q x + g^\top x + c - \kappa\bm{1}^T \log(d) \\
\st \quad & Ax + b = 0 \\
          & Gx + h = d \\
          & Lx + l = b_\kappa(\sigma) \\
          & Rx + r = b_\kappa(-\sigma)
\end{aligned}
\end{equation}
We then introduce the Augmented Lagrangian for \eqref{qpcc_relax} as:
\begin{equation}\label{aug_lag}
\begin{aligned}
\mathcal{L}_\rho(x, \sigma, d; \alpha, \beta, \gamma_L, \gamma_R) = & \frac12 x^\top Q x + g^\top x + c - \kappa\bm{1}^T \log(d) \\
& + \alpha^\top (Ax + b) + \frac\rho2 \norm{Ax + b}^2 \\
& + \beta^\top (Gx + h - d) + \frac\rho2 \norm{Gx+h-d}^2 \\
& + \gamma_L^\top (Lx + l - b_\kappa(\sigma)) + \frac\rho2 \norm{Lx + l - b_\kappa(\sigma)}^2 \\
& + \gamma_R^\top (Rx + r - b_\kappa(-\sigma)) + \frac\rho2 \norm{Rx + r - b_\kappa(-\sigma)}^2
\end{aligned}
\end{equation}
Using an Augmented-Lagrangian approach, we can solve \eqref{qpcc_relax} by iteratively minimizing \eqref{aug_lag} 
over primal variables $x, \sigma, d$ and updating dual variables $\alpha, \beta, \gamma_L, \gamma_R$.



We first derive stationarity with respect to the original primal variable $x$. Define the following 
shorthands:
\begin{align}
    r_A &:= Ax + b \\
    r_G &:= Gx + h - d \\
    r_L &:= Lx + l - b_\kappa(\sigma) \\
    r_R &:= Rx + r - b_\kappa(-\sigma)
\end{align}
We therefore have the following stationarity condition:
\begin{align}
\pd{\mathcal{L}_\rho}{x} &= Qx + g + A^\top \alpha + \rho A^\top r_A + G^\top \beta + \rho G^\top r_G + L^\top \gamma_L + \rho L^\top r_L + R^\top \gamma_R + \rho R^\top r_R = 0 \\
&= Qx + g + A^\top \lambda + G^\top \mu + L^\top \tau_L + R^\top \tau_R = 0
\end{align}
We introduce the primal-dual shifted multiplier variables:
\begin{align}
    \lambda &= \alpha + \rho r_A = \alpha + \rho (Ax + b) \\
    \mu &= \beta + \rho r_G = \beta + \rho (Gx + h - d) \\
    \tau_L &= \gamma_L + \rho r_L = \gamma_L + \rho (Lx + l - b_\kappa(\sigma)) \\
    \tau_R &= \gamma_R + \rho r_R = \gamma_R + \rho (Rx + r - b_\kappa(-\sigma))
\end{align}

For variable $d$, we have the following stationarity condition:
\begin{align}
\pd{\mathcal{L}_\rho}{d} &= -\beta - \rho r_G - \kappa d^{-1} = 0 \\
&= -\mu - \kappa d^{-1} = 0 \Longleftrightarrow \mu \odot d = -\kappa \bm{1}
\end{align}
We can satisfy this relaxed complementarity condition by again using \eqref{retraction_map} as:
\begin{equation}
    d = b_\kappa(v), \quad \mu = -b_\kappa(-v), \quad v \in \R^p
\end{equation}

And finally for variable $\sigma$, we have the following stationarity condition:
\begin{align}
\pd{\mathcal{L}_\rho}{\sigma} &= -b'_\kappa(\sigma) (\gamma_L + \rho r_L) + b'_\kappa(-\sigma) (\gamma_R + \rho r_R) = 0 \\
&= -b'_\kappa(\sigma) \tau_L + b'_\kappa(-\sigma) \tau_R = 0
\end{align}

In all, we have the following KKT conditions for \eqref{qpcc_relax}. We note that
$b''_\kappa(\sigma) = b''_\kappa(-\sigma)$ since $b''_\kappa(x)$ is an even function.
\begin{align}
    r_x &:= Qx + g + A^\top \lambda + G^\top \mu + L^\top \tau_L + R^\top \tau_R &&= 0 \\
    r_v &:= \mu + b_\kappa(-v) &&= 0 \\
    r_\sigma &:= -b'_\kappa(\sigma) \tau_L + b'_\kappa(-\sigma) \tau_R &&= 0 \\
    r_\lambda &:= Ax + b - \rho^{-1} (\lambda - \alpha) &&= 0 \\
    r_\mu &:= Gx + h - b_\kappa(v) - \rho^{-1} (\mu - \beta) &&= 0 \\
    r_{\tau_L} &:= Lx + l - b_\kappa(\sigma) - \rho^{-1} (\tau_L - \gamma_L) &&= 0 \\
    r_{\tau_R} &:= Rx + r - b_\kappa(-\sigma) - \rho^{-1} (\tau_R - \gamma_R) &&= 0
\end{align}
To solve this root-finding problem, we first write the asymmetric Newton system as the following,
where we note $b''_\kappa(\sigma) = b''_\kappa(-\sigma)$:
\begin{equation}
\begin{bmatrix}
    Q & 0 & 0 & A^\top & G^\top & L^\top & R^\top \\
    0 & -b'_\kappa(-v) & 0 & 0 & I & 0 & 0 \\
    0 & 0 & -b''_\kappa(\sigma) \odot (\tau_L + \tau_R) & 0 & 0 & -b'_\kappa(\sigma) & b'_\kappa(-\sigma) \\
    A & 0 & 0 & -\rho^{-1} I & 0 & 0 & 0 \\
    G & -b'_\kappa(v) & 0 & 0 & -\rho^{-1} I & 0 & 0 \\
    L & 0 & -b'_\kappa(\sigma) & 0 & 0 & -\rho^{-1} I & 0 \\
    R & 0 & b'_\kappa(-\sigma) & 0 & 0 & 0 & -\rho^{-1} I
\end{bmatrix}
\begin{bmatrix}
    \Delta x \\ \Delta v \\ \Delta \sigma \\ \Delta \lambda \\ \Delta \mu \\ \Delta \tau_L \\ \Delta \tau_R
\end{bmatrix}
=
-\begin{bmatrix}
    r_x \\ r_v \\ r_\sigma \\ r_\lambda \\ r_\mu \\ r_{\tau_L} \\ r_{\tau_R}
\end{bmatrix}
\end{equation}
Using the second equation $-b'_\kappa(-v) \Delta v + \Delta \mu = -r_v$, we solve for $\Delta \mu$ and get:
\[
  Q \Delta x + A^\top \Delta \lambda + G^\top b'_\kappa(-v) \Delta v + L^\top \Delta \tau_L + R^\top \Delta \tau_R = -r_x + G^\top r_v
\]
We then subtract $G^\top$ times the fifth equation $G \Delta x - b'_\kappa(v) \Delta v - \rho^{-1} \Delta \mu = -r_\mu$ to get:
\[
    (Q - G^\top G) \Delta x + A^\top \Delta \lambda + G^\top(I + \rho^{-1}b_\kappa'(-v)) \Delta v + L^\top \Delta \tau_L + R^\top \Delta \tau_R = -r_x + G^\top (r_v + r_\mu + \rho^{-1} r_v)
\]
% -1/rho \delta mu =
% -1/rho (b_k'(-v) \delta v - r_v ) 

\begin{equation}
\begin{bmatrix}
    Q - G^\top G & G^\top (I + \rho^{-1}b_\kappa'(-v)) & 0 & A^\top & G^\top & L^\top & R^\top \\
    G & -b'_\kappa(v) & 0 & 0 & -\rho^{-1} I & 0 & 0 \\
    % 0 & -b'_\kappa(-v) & 0 & 0 & I & 0 & 0 \\
    0 & 0 & -b''_\kappa(\sigma) \odot (\tau_L + \tau_R) & 0 & 0 & -b'_\kappa(\sigma) & b'_\kappa(-\sigma) \\
    A & 0 & 0 & -\rho^{-1} I & 0 & 0 & 0 \\
    L & 0 & -b'_\kappa(\sigma) & 0 & 0 & -\rho^{-1} I & 0 \\
    R & 0 & b'_\kappa(-\sigma) & 0 & 0 & 0 & -\rho^{-1} I
\end{bmatrix}
\end{equation}

\newpage
We symmetrize this system by multiplying the second set of equations by $-b'_\kappa(v)$
to arrive at the following symmetric Newton system:
\begin{equation}
\underbrace{
\begin{bmatrix}
    Q & 0 & 0 & A^\top & G^\top & L^\top & R^\top \\
    0 & b'_\kappa(v) b'_\kappa(-v) & 0 & 0 & -b'_\kappa(v) & 0 & 0 \\
    0 & 0 & -b''_\kappa(\sigma) \odot (\tau_L + \tau_R) & 0 & 0 & -b'_\kappa(\sigma) & b'_\kappa(-\sigma) \\
    A & 0 & 0 & -\rho^{-1} I & 0 & 0 & 0 \\
    G & -b'_\kappa(v) & 0 & 0 & -\rho^{-1} I & 0 & 0 \\
    L & 0 & -b'_\kappa(\sigma) & 0 & 0 & -\rho^{-1} I & 0 \\
    R & 0 & b'_\kappa(-\sigma) & 0 & 0 & 0 & -\rho^{-1} I
\end{bmatrix}
}_K
\begin{bmatrix}
    \Delta x \\ \Delta v \\ \Delta \sigma \\ \Delta \lambda \\ \Delta \mu \\ \Delta \tau_L \\ \Delta \tau_R
\end{bmatrix}
=
-\begin{bmatrix}
    r_x \\ -b'_\kappa(v) r_v \\ r_\sigma \\ r_\lambda \\ r_\mu \\ r_{\tau_L} \\ r_{\tau_R}
\end{bmatrix}
\end{equation}

We can decompose this symmetric KKT system into the following block form:
\begin{equation}
K =
\begin{bmatrix}
    H & J^\top \\
    J & -\rho^{-1} I
\end{bmatrix}
\end{equation}
where:
\begin{equation}
H = 
\begin{bmatrix}
    Q & 0 & 0  \\
    0 & b'_\kappa(v) b'_\kappa(-v) & 0 \\
    0 & 0 & -b''_\kappa(\sigma) \odot (\tau_L + \tau_R)
\end{bmatrix},
\qquad
J = 
\begin{bmatrix}
    A & 0 & 0 \\
    G & -b'_\kappa(v) & 0 \\
    L & 0 & -b'_\kappa(\sigma) \\
    R & 0 & b'_\kappa(-\sigma)
\end{bmatrix}
\end{equation}

\subsection{Inertia Correction}
By Haynsworth's inertia additivity, the inertia of the KKT matrix is given by:
\begin{equation}
    \inertia(K) = \inertia(-\rho^{-1} I) + \inertia(H + \rho J^\top J)
\end{equation}
The target inertia for the KKT system is $\inertia(K) = (n + p + c, m + p + 2c, 0)$;
since $\inertia(-\rho^{-1} I) = (0, m + p + 2c, 0)$, we require that $\inertia(H + \rho J^\top J) = (n + p + c, 0, 0)$,
or equivalently that $H + \rho J^\top J \succ 0$. We can ensure this by adding a diagonal regularization term $\delta I$ to $H$ such that:
\begin{equation}
    H + \delta I + \rho J^\top J \succ 0
\end{equation}

For further analysis, we write out the matrix $\rho J^\top J$ as:
\begin{equation}
\rho J^\top J = \rho
\begin{bmatrix}
    A^\top A + G^\top G + L^\top L + R^\top R & -G^\top b'_\kappa(v) & -L^\top b'_\kappa(\sigma) + R^\top b'_\kappa(-\sigma) \\
    -b'_\kappa(v) G & b'_\kappa(v)^2 & 0 \\
    -b'_\kappa(\sigma) L + b'_\kappa(-\sigma) R & 0 & b'_\kappa(\sigma)^2 + b'_\kappa(-\sigma)^2
\end{bmatrix}
\end{equation}

% -B_k(-v) dv + dmu = -r_v 
% dv = B_k(-v)^{-1}(r_v + dmu)

% G*dx + -B_k(v) dv - 1/p * d mu = -r_mu
% G*dx + -B_k(v) * B_k(-v)^{-1}(r_v + dmu)
% G*dx + 

\section{Algorithm}

\end{document}